\chapter{Theta modularity, lattice dissection, and sums of squares}
\label{chap:theta-modularity-squares}

This chapter collects the theta-function consequences that are not needed for
the elementary development of the Jacobi triple product in
\Cref{chap:jtp}, but which reveal its modular and arithmetic force.  The
arguments are deliberately explicit: Poisson summation supplies modular
inversion, lattice dissection supplies product identities, and unique
factorization or logarithmic differentiation converts theta coefficients into
divisor sums.  Throughout the analytic sections, $0<\AbsoluteValue q<1$.

\section{Watson's quintuple product}

\begin{theorem}[Watson's quintuple product]
\label{qg:thm-quintuple-product}
For $z\ne0$,
\begin{align}
 &(q,-z,-q/z;q)_\infty(qz^2,q/z^2;q^2)_\infty\notag\\
 &\qquad=\sum_{n\in\IntegerNumbers}(-1)^n
 q^{n(3n-1)/2}z^{3n}(1+zq^n).
\label{eq:qg-quintuple-product}
\end{align}
\end{theorem}

\begin{proof}
By \Cref{thm:jtp},
\begin{align*}
 (-z,-q/z,q;q)_\infty
   &=\sum_{r\in\IntegerNumbers}q^{r(r-1)/2}z^r,\\
 (qz^2,q/z^2,q^2;q^2)_\infty
   &=\sum_{s\in\IntegerNumbers}(-1)^sq^{s^2}z^{2s}.
\end{align*}
Consequently, the left side of
\eqref{eq:qg-quintuple-product}, multiplied by $(q^2;q^2)_\infty$, is
the product of these two bilateral series.  Both converge absolutely for $0<\AbsoluteValue q<1$
and $z\ne0$, so the doubly indexed family of products is unconditionally summable over
$\IntegerNumbers^2$; grouping it along the fibres $r+2s=m$ is therefore a legitimate
rearrangement, and no appeal to uniqueness of Laurent coefficients is required.  The coefficient
of $z^m$ is
\begin{equation}
 q^{m(m-1)/2}\sum_{s\in\IntegerNumbers}(-1)^s
 q^{3s^2+(1-2m)s}.
\label{eq:qg-quintuple-coefficient}
\end{equation}
Apply the minus-sign form of the triple product with base $Q=q^6$, together with the iterated
quasi-periodicity
\[
 \BilateralTheta{Q^{j}w}{Q}\,Q^{\binom j2}w^{\,j}=\BilateralTheta{w}{Q}
 \qquad(j\in\IntegerNumbers),
\]
which follows from \Cref{cor:theta-quasi} by induction in both directions.  This gives, when
$m=3n$,
\[
 \sum_{s\in\IntegerNumbers}(-1)^sq^{3s^2+(1-6n)s}
 =(-1)^nq^{-3n^2+n}(q^2;q^2)_\infty.
\]
Thus \eqref{eq:qg-quintuple-coefficient} equals
$(-1)^nq^{n(3n-1)/2}(q^2;q^2)_\infty$.  When $m=3n+1$, the same
calculation gives
$(-1)^nq^{n(3n+1)/2}(q^2;q^2)_\infty$; explicitly, the shift is by $Q^{-n}$ starting from $-q^{2}$
rather than from $-q^{4}$, which contributes $(-1)^nq^{-3n^2-n}$, and
$\binom{3n+1}{2}-3n^2-n=n(3n+1)/2$.

When $m=3n+2$ the inner sum is $\BilateralTheta{-q^{-6n}}{Q}=\BilateralTheta{-Q^{-n}}{Q}$, and
$-Q^{-n}$ lies in the zero set $\{-Q^{j}:j\in\IntegerNumbers\}$ of $\BilateralTheta{\cdot}{Q}$
given by \Cref{cor:theta-quasi}; hence it vanishes, uniformly in $n$ and with no case
distinction.  The shift used above is equally uniform: taking $w=-1$ and $j=-n$ in the iterated
quasi-periodicity gives $\BilateralTheta{-Q^{-n}}{Q}\,Q^{\binom{-n}{2}}(-1)^{-n}
=\BilateralTheta{-1}{Q}=0$, since $-1=-Q^{0}$ lies in the zero set.  What does require a case
distinction is reading the vanishing directly off the product without shifting: the offending
factor then sits in $(q^{-6n};q^6)_\infty$ when $n\ge0$, but in $(q^{6n+6};q^6)_\infty$ when
$n\le-1$.

After division by the nonzero product $(q^2;q^2)_\infty$, only the powers
$z^{3n}$ and $z^{3n+1}$ remain.  Their coefficients are respectively
\[
 (-1)^nq^{n(3n-1)/2},
 \qquad
 (-1)^nq^{n(3n+1)/2},
\]
which combine to the right side of
\eqref{eq:qg-quintuple-product}.
\end{proof}

\begin{remark}[Coefficient dissection and root-system products]
The proof illustrates a mechanism that persists in higher-rank product
identities.  Multiplying two Jacobi factors first produces a lattice sum;
fixing the exponent of $z$ restricts that lattice to an affine slice; and the
remaining one-dimensional theta factor then uses its zeros and
quasi-periodicity to distinguish residue classes.  In Macdonald-type
root-system products, higher-dimensional lattice cosets and Weyl symmetries
play analogous roles.  This is a structural comparison, not an invocation or
proof of a root-system identity.
\end{remark}

\section{Jacobi theta functions and the heat equation}

Let $\tau$ lie in the upper half-plane and put
\[
 q=\EulerE^{\pi\ImaginaryUnit\tau},
 \qquad
 \zeta=\EulerE^{2\pi\ImaginaryUnit z}.
\]
Whenever a rational power of $q$ occurs in the theta or eta formulas, it is
defined by
\[
 q^\alpha:=\EulerE^{\pi\ImaginaryUnit\alpha\tau};
\]
it is not a power obtained from an unspecified logarithm of the complex number
$q$.
Whenever $\zeta^{1/2}$ occurs below, it means the single-valued expression
$\EulerE^{\pi\ImaginaryUnit z}$, not an unspecified square root.

\begin{definition}[Jacobi theta functions]
\label{qg:def-jacobi-theta}
\begin{align}
 \vartheta_3(z\mid\tau)
  &=\sum_{n\in\IntegerNumbers}
    \EulerE^{\pi\ImaginaryUnit n^2\tau+2\pi\ImaginaryUnit nz}
   =\sum_{n\in\IntegerNumbers}q^{n^2}\zeta^n,
 \label{eq:qg-theta3-definition}\\
 \vartheta_4(z\mid\tau)
  &=\sum_{n\in\IntegerNumbers}(-1)^n
    \EulerE^{\pi\ImaginaryUnit n^2\tau+2\pi\ImaginaryUnit nz},
 \label{eq:qg-theta4-definition}\\
 \vartheta_2(z\mid\tau)
  &=\sum_{n\in\IntegerNumbers}
    \EulerE^{\pi\ImaginaryUnit(n+1/2)^2\tau
      +2\pi\ImaginaryUnit(n+1/2)z}.
 \label{eq:qg-theta2-definition}
\end{align}
\end{definition}

These series and all their termwise derivatives converge normally on compact
subsets of $\ComplexNumbers\times\{\operatorname{Im}\tau>0\}$.

\begin{proposition}[Theta product forms]
\label{qg:prop-theta-products}
One has
\begin{align}
 \vartheta_3(z\mid\tau)
  &=(q^2,-q\zeta,-q/\zeta;q^2)_\infty,
 \label{eq:qg-theta3-product}\\
 \vartheta_4(z\mid\tau)
  &=(q^2,q\zeta,q/\zeta;q^2)_\infty,
 \label{eq:qg-theta4-product}\\
 \vartheta_2(z\mid\tau)
  &=q^{1/4}\zeta^{1/2}
    (q^2,-q^2\zeta,-\zeta^{-1};q^2)_\infty.
 \label{eq:qg-theta2-product}
\end{align}
In particular,
\begin{align}
 \vartheta_2(0\mid\tau)
  &=2q^{1/4}(q^2;q^2)_\infty(-q^2;q^2)_\infty^2,
 \label{eq:qg-theta2-null}\\
 \vartheta_3(0\mid\tau)
  &=(q^2;q^2)_\infty(-q;q^2)_\infty^2,
 \label{eq:qg-theta3-null}\\
 \vartheta_4(0\mid\tau)
  &=(q^2;q^2)_\infty(q;q^2)_\infty^2.
 \label{eq:qg-theta4-null}
\end{align}
\end{proposition}

\begin{proof}
Apply \Cref{thm:jtp} with base $q^2$.  The substitution $z=-q\zeta$
there works because
\[
 (-1)^n(q^2)^{\binom n2}(-q\zeta)^n=q^{n^2}\zeta^n;
\]
the substitution $z=q\zeta$ gives the second formula.  Finally,
\[
 \vartheta_2(z\mid\tau)
 =q^{1/4}\zeta^{1/2}
  \sum_{n\in\IntegerNumbers}q^{n(n+1)}\zeta^n,
\]
and the triple product with argument $-q^2\zeta$ gives
\eqref{eq:qg-theta2-product}.  Set $z=0$ and use
$(-1;q^2)_\infty=2(-q^2;q^2)_\infty$ to obtain the three null values.
\end{proof}

\begin{proposition}[Heat equation]
\label{qg:prop-theta-heat}
Each $\vartheta_j$, for $j=2,3,4$, satisfies
\begin{equation}
 \frac{\partial\vartheta_j}{\partial\tau}
 =\frac{1}{4\pi\ImaginaryUnit}
   \frac{\partial^2\vartheta_j}{\partial z^2}.
\label{eq:qg-theta-heat}
\end{equation}
\end{proposition}

\begin{proof}
For a summand with shift $\alpha\in\{0,1/2\}$, differentiation in
$\tau$ multiplies by $\pi\ImaginaryUnit(n+\alpha)^2$, whereas two
differentiations in $z$ multiply by $-4\pi^2(n+\alpha)^2$.  Since
$-4\pi^2/(4\pi\ImaginaryUnit)=\pi\ImaginaryUnit$, the identities hold
term by term; normal convergence justifies both differentiations.
\end{proof}

\section{Poisson summation and modular inversion}

\begin{definition}[Fourier transform convention]
\label{qg:def-fourier-transform}
For an integrable function $f:\RealNumbers\to\ComplexNumbers$ and
$\xi\in\RealNumbers$, write
\[
 \FourierTwoPi{f}(\xi)
 =\int_{\RealNumbers}f(x)\EulerE^{-2\pi\ImaginaryUnit x\xi}
   \,\Differential x.
\]
\end{definition}

\begin{theorem}[Poisson summation]
\label{qg:thm-poisson}
For every Schwartz function $f:\RealNumbers\to\ComplexNumbers$, one has
\begin{equation}
 \sum_{n\in\IntegerNumbers}f(n)
 =\sum_{k\in\IntegerNumbers}\FourierTwoPi{f}(k).
\label{eq:qg-poisson}
\end{equation}
\end{theorem}

\begin{proof}
The periodization $P(x)=\sum_{n\in\IntegerNumbers}f(x+n)$ and all its
derivatives converge uniformly.  Thus $P$ is smooth and $1$-periodic.  Its
$k$th Fourier coefficient is
\begin{align*}
 \int_0^1P(x)\EulerE^{-2\pi\ImaginaryUnit kx}\,\Differential x
 &=\sum_{n\in\IntegerNumbers}\int_0^1
   f(x+n)\EulerE^{-2\pi\ImaginaryUnit kx}\,\Differential x\\
 &=\int_{\RealNumbers}f(u)\EulerE^{-2\pi\ImaginaryUnit ku}
   \,\Differential u
 =\FourierTwoPi{f}(k),
\end{align*}
because $\EulerE^{2\pi\ImaginaryUnit kn}=1$.  Repeated integration by
parts makes these Fourier coefficients rapidly decreasing.  The Fourier
inversion theorem for smooth periodic functions therefore gives an absolutely
and uniformly convergent Fourier series for $P$.  Evaluating it at $x=0$ proves
\eqref{eq:qg-poisson}.
\end{proof}

\begin{lemma}[Fourier transform of a complex Gaussian]
\label{qg:lem-gaussian-fourier}
If $\operatorname{Im}\tau>0$ and $u\in\ComplexNumbers$, then, with the
principal square-root branch,
\begin{equation}
 \int_{\RealNumbers}
 \EulerE^{\pi\ImaginaryUnit\tau x^2+2\pi\ImaginaryUnit ux}
 \,\Differential x
 =(-\ImaginaryUnit\tau)^{-1/2}
  \EulerE^{-\pi\ImaginaryUnit u^2/\tau}.
\label{eq:qg-complex-gaussian}
\end{equation}
\end{lemma}

\begin{proof}
First let $\tau=\ImaginaryUnit t$ with $t>0$, and for real $u$ put
\[
 I_t(u)=\int_{\RealNumbers}
 \EulerE^{-\pi t x^2+2\pi\ImaginaryUnit ux}\,\Differential x.
\]
Differentiation under the integral is justified by Gaussian domination.  Since
\[
 \frac{\Differential}{\Differential x}\EulerE^{-\pi t x^2}
 =-2\pi t x\,\EulerE^{-\pi t x^2},
\]
integration by parts, with vanishing boundary term, gives
\begin{align*}
 I_t'(u)
 &=2\pi\ImaginaryUnit\int_{\RealNumbers}
   x\EulerE^{-\pi t x^2+2\pi\ImaginaryUnit ux}\,\Differential x\\
 &=-\frac{\ImaginaryUnit}{t}
   \int_{\RealNumbers}
   \frac{\Differential}{\Differential x}
   \bigl(\EulerE^{-\pi t x^2}\bigr)
   \EulerE^{2\pi\ImaginaryUnit ux}\,\Differential x\\
 &=-\frac{2\pi u}{t}I_t(u).
\end{align*}
Because the ordinary Gaussian integral gives $I_t(0)=t^{-1/2}$, this
differential equation yields
\[
 I_t(u)=t^{-1/2}\EulerE^{-\pi u^2/t}
 \qquad (u\in\RealNumbers).
\]

For $u$ in a compact subset of $\ComplexNumbers$, the integrand and all its
$u$-derivatives are dominated by a constant times
\[
 \EulerE^{-\pi t x^2+C\AbsoluteValue x},
\]
so the integral is entire in $u$.  The identity theorem therefore extends the
formula to every complex $u$.

Finally, if $K$ is a compact subset of the upper half-plane and $u$ is fixed,
then for some $\delta>0$ and all $\tau\in K$,
\[
 \left|
 \EulerE^{\pi\ImaginaryUnit\tau x^2+2\pi\ImaginaryUnit ux}
 \right|
 \le
 \EulerE^{-\pi\delta x^2+2\pi\AbsoluteValue{\operatorname{Im}u}
                    \AbsoluteValue x}.
\]
The same bound with an additional polynomial factor handles every
$\tau$-derivative.  Thus the integral is holomorphic in $\tau$.  The identity
theorem extends the formula from the positive imaginary axis to the whole
upper half-plane.  Since $-\ImaginaryUnit\tau$ lies in the right half-plane,
its principal square root is holomorphic there and is characterized by
\[
 (-\ImaginaryUnit\cdot\ImaginaryUnit t)^{1/2}=t^{1/2}.
\]
\end{proof}

\begin{theorem}[Jacobi imaginary transformation]
\label{qg:thm-jacobi-imaginary}
For $\operatorname{Im}\tau>0$ and $z\in\ComplexNumbers$,
\begin{equation}
 \vartheta_3\!\left(\frac z\tau\middle|-\frac1\tau\right)
 =(-\ImaginaryUnit\tau)^{1/2}
  \EulerE^{\pi\ImaginaryUnit z^2/\tau}\vartheta_3(z\mid\tau).
\label{eq:qg-theta3-modular}
\end{equation}
At $z=0$,
\begin{align}
 \vartheta_2(0\mid-1/\tau)
  &=(-\ImaginaryUnit\tau)^{1/2}\vartheta_4(0\mid\tau),
 \label{eq:qg-theta2-modular}\\
 \vartheta_3(0\mid-1/\tau)
  &=(-\ImaginaryUnit\tau)^{1/2}\vartheta_3(0\mid\tau),
 \label{eq:qg-theta3-null-modular}\\
 \vartheta_4(0\mid-1/\tau)
  &=(-\ImaginaryUnit\tau)^{1/2}\vartheta_2(0\mid\tau).
 \label{eq:qg-theta4-modular}
\end{align}
\end{theorem}

\begin{proof}
Apply \Cref{qg:thm-poisson} to
$f(x)=\EulerE^{\pi\ImaginaryUnit\tau x^2+2\pi\ImaginaryUnit zx}$.
For fixed $z$, the Gaussian decay supplied by $\operatorname{Im}\tau>0$
shows that $f$ is a Schwartz function.
By \Cref{qg:lem-gaussian-fourier},
\[
 \FourierTwoPi{f}(k)=(-\ImaginaryUnit\tau)^{-1/2}
 \EulerE^{-\pi\ImaginaryUnit(z-k)^2/\tau}.
\]
After factoring out
$\EulerE^{-\pi\ImaginaryUnit z^2/\tau}$, Poisson summation is exactly
\eqref{eq:qg-theta3-modular}.  For the shifted Gaussian
$f(x)=\EulerE^{\pi\ImaginaryUnit\tau(x+1/2)^2}$, translation by $1/2$
multiplies the $k$th Fourier transform by $(-1)^k$.  Poisson summation then
gives
\[
 \vartheta_2(0\mid\tau)
 =(-\ImaginaryUnit\tau)^{-1/2}\vartheta_4(0\mid-1/\tau).
\]
This is \eqref{eq:qg-theta4-modular}; replacing $\tau$ by $-1/\tau$
and using the chosen branch gives \eqref{eq:qg-theta2-modular}.
\end{proof}

\begin{definition}[Dedekind eta function]
\label{qg:def-dedekind-eta}
For $\operatorname{Im}\tau>0$,
\begin{equation}
 \eta(\tau)=\EulerE^{\pi\ImaginaryUnit\tau/12}
 \prod_{n=1}^{\infty}(1-\EulerE^{2\pi\ImaginaryUnit n\tau})
 =q^{1/12}(q^2;q^2)_\infty.
\label{eq:qg-eta-definition}
\end{equation}
\end{definition}

\begin{theorem}[Theta--eta product and eta modularity]
\label{qg:thm-theta-eta-modular}
One has
\begin{align}
 \vartheta_2(0\mid\tau)\vartheta_3(0\mid\tau)
 \vartheta_4(0\mid\tau)&=2\eta(\tau)^3,
 \label{eq:qg-theta-eta}\\
 \eta(\tau+1)&=\EulerE^{\pi\ImaginaryUnit/12}\eta(\tau),
 \label{eq:qg-eta-T}\\
 \eta(-1/\tau)&=(-\ImaginaryUnit\tau)^{1/2}\eta(\tau).
 \label{eq:qg-eta-S}
\end{align}
\end{theorem}

\begin{proof}
Multiply the three null-value formulas in
\Cref{qg:prop-theta-products}.  The elementary product relations
\[
 (-q;q^2)_\infty(q;q^2)_\infty=(q^2;q^4)_\infty,
 \qquad
 (-q^2;q^2)_\infty=\frac{(q^4;q^4)_\infty}{(q^2;q^2)_\infty},
\]
and
$(q^2;q^2)_\infty=(q^2;q^4)_\infty(q^4;q^4)_\infty$
reduce the product to $2q^{1/4}(q^2;q^2)_\infty^3=2\eta(\tau)^3$.

The translation formula follows immediately from the definition.  Applying
the three modular formulas in \Cref{qg:thm-jacobi-imaginary} to the theta--eta
product yields
\[
 \eta(-1/\tau)^3=(-\ImaginaryUnit\tau)^{3/2}\eta(\tau)^3.
\]
Moreover, $\eta$ has no zeros in the upper half-plane.  Indeed,
\[
 \sum_{n\ge1}\AbsoluteValue{\EulerE^{2\pi\ImaginaryUnit n\tau}}<\infty,
\]
and none of the factors
$1-\EulerE^{2\pi\ImaginaryUnit n\tau}$ vanishes; hence the convergent infinite
product in \eqref{eq:qg-eta-definition} is nonzero.  Since
$-\ImaginaryUnit\tau$ lies in the right half-plane, its principal square root
is also holomorphic and nonzero there.
The quotient
\[
 h(\tau)=\frac{\eta(-1/\tau)}
 {(-\ImaginaryUnit\tau)^{1/2}\eta(\tau)}
\]
is holomorphic and takes values among the cube roots of unity, hence is
constant.  At $\tau=\ImaginaryUnit$ it equals $1$, proving
\eqref{eq:qg-eta-S} with no unresolved root of unity.
\end{proof}

\section{Ramanujan's theta function and lattice dissection}

\begin{definition}[Ramanujan's general theta function]
\label{qg:def-ramanujan-theta}
For $a,b\in\ComplexNumbers$ with $\AbsoluteValue{ab}<1$, set
\begin{equation}
 \RamanujanTheta{a}{b}=\sum_{n\in\IntegerNumbers}
 a^{n(n+1)/2}b^{n(n-1)/2}.
\label{eq:qg-ramanujan-theta}
\end{equation}
Here a zeroth power is $1$, including at a zero base.  Pairing the terms
indexed by $n$ and $-n$ gives the unambiguous form
\[
 \RamanujanTheta{a}{b}=1+\sum_{n=1}^{\infty}
 (a^n+b^n)(ab)^{n(n-1)/2}.
\]
This series converges normally on compact subsets of
$\{(a,b):\AbsoluteValue{ab}<1\}$: on such a compact set its $n$th paired
term is bounded by $2M^n\rho^{n(n-1)/2}$ for suitable $M<\infty$ and
$\rho<1$.
\end{definition}

\begin{proposition}[Ramanujan's product]
\label{qg:prop-ramanujan-product}
For $a,b\in\ComplexNumbers$ with $\AbsoluteValue{ab}<1$,
\begin{equation}
 \RamanujanTheta{a}{b}=(-a,-b,ab;ab)_\infty.
\label{eq:qg-ramanujan-product}
\end{equation}
At base zero, the infinite symbol is interpreted by its natural extension
$(c;0)_\infty:=1-c$.
\end{proposition}

\begin{proof}
If $ab\ne0$, then both $a$ and $b$ are nonzero.  In \Cref{thm:jtp}, take
the base to be $ab$ and the argument to be $-a$.  The $n$th summand becomes
\[
 (-1)^n(ab)^{\binom n2}(-a)^n
 =a^{n(n+1)/2}b^{n(n-1)/2},
\]
and the product is $(-a,-b,ab;ab)_\infty$.

If $ab=0$, the paired form of the defining series gives
$\RamanujanTheta{a}{b}=1+a+b$.  With the stated zero-base convention, the product side is
\[
 (1+a)(1+b)=1+a+b+ab=1+a+b.
\]
Thus the identity also holds in the degenerate cases.
\end{proof}

For $\AbsoluteValue Q<1$ and $x\in\ComplexNumbers^\times$ write
\[
 \BilateralQuadraticTheta{x}{Q}
 =\sum_{n\in\IntegerNumbers}Q^{n^2}x^n.
\]
The quadratic exponent makes this series absolutely convergent for every such
$Q$ and $x$.
Its exact crosswalk to the bilateral Jacobi series defined earlier is
\[
 \BilateralQuadraticTheta{x}{Q}=\BilateralTheta{Qx}{Q^2}.
\]
The semantic subscript distinguishes this function from the Fabius transfer
operator $\mathcal T$ and from every proof-local transformation.
With the theta nome and elliptic exponential declared above, the complete
crosswalk is
\begin{align*}
 \vartheta_3(z\mid\tau)
 &=\BilateralQuadraticTheta{\zeta}{q}
  =\BilateralTheta{q\zeta}{q^2},\\
 \vartheta_4(z\mid\tau)
 &=\BilateralQuadraticTheta{-\zeta}{q}
  =\MultiplicativeTheta{q\zeta}{q^2},\\
 \vartheta_2(z\mid\tau)
 &=q^{1/4}\zeta^{1/2}\BilateralTheta{q^2\zeta}{q^2}.
\end{align*}

\begin{theorem}[Schr\"oter lattice dissection]
\label{qg:thm-schroeter}
Let $a,b$ be positive integers, let $0<\AbsoluteValue q<1$, and let
$x,y\in\ComplexNumbers^\times$.  Then
\begin{align}
 \BilateralQuadraticTheta{x}{q^a}\BilateralQuadraticTheta{y}{q^b}
 &=\sum_{k=0}^{a+b-1}y^kq^{bk^2}
 \BilateralQuadraticTheta{xyq^{2bk}}{q^{a+b}}\notag\\
 &\qquad\qquad\times
 \BilateralQuadraticTheta{y^ax^{-b}q^{2abk}}{q^{ab(a+b)}}.
\label{eq:qg-schroeter}
\end{align}
\end{theorem}

\begin{proof}
Expand the left side over $(n,m)\in\IntegerNumbers^2$ and partition the
lattice by the unique $k\in\{0,\ldots,a+b-1\}$ for which
$m-n=k+(a+b)s$ with $s\in\IntegerNumbers$.  In that residue class the
change of variables
\[
 n=r-bs,
 \qquad
 m=k+r+as
\]
is a bijection with $(r,s)\in\IntegerNumbers^2$.  Direct substitution gives
\begin{align*}
 x^ny^m&=y^k(xy)^r(y^ax^{-b})^s,\\
 an^2+bm^2
 &=(a+b)r^2+2bkr+bk^2
   +ab(a+b)s^2+2abks.
\end{align*}
The $r$- and $s$-sums are therefore the two displayed theta series, and the
remaining factor is $y^kq^{bk^2}$.  Absolute convergence justifies the
partition and rearrangement.
\end{proof}

\begin{remark}[Sublattice cosets and diagonalization]
For fixed $k$, the linear part of the change of variables is
\[
 \binom{n}{m-k}
 =\begin{pmatrix}1&-b\\1&a\end{pmatrix}\binom{r}{s},
\]
whose determinant is $a+b$.  Its image is the index-$(a+b)$ sublattice on
which the second coordinate minus the first is divisible by $a+b$.
Translating this image by $(0,k)$ gives the coset
$m-n\equiv k\pmod{a+b}$, and the values $k=0,\ldots,a+b-1$ give all cosets.
Equivalently, this is the rank-two Smith-normal-form description of the
congruence map $(n,m)\mapsto m-n\pmod{a+b}$.  At the same time, the mixed
$rs$ term in
$an^2+bm^2$ cancels under this basis change.  The remaining linear terms,
which depend on the coset representative $k$, are absorbed into the arguments
of the two theta series.  Thus Schr\"oter dissection is simultaneously a
sublattice decomposition and a diagonalization of the quadratic form; this
remark records the mechanism and asserts no additional identity.
\end{remark}

\section{Theta powers and sums of squares}

Let $r_d(n)$ be the number of ordered $d$-tuples
$(x_1,\ldots,x_d)\in\IntegerNumbers^d$ with
$x_1^2+\cdots+x_d^2=n$.
For $d=0$ we use the unique empty tuple, so
$r_0(0)=1$ and $r_0(n)=0$ for $n\ge1$.

\begin{proposition}[Universal sums-of-squares generating function]
\label{qg:prop-squares-theta}
For every nonnegative integer $d$ and $\AbsoluteValue q<1$,
\begin{equation}
 \left(\sum_{m\in\IntegerNumbers}q^{m^2}\right)^d
 =\sum_{n=0}^{\infty}r_d(n)q^n.
\label{eq:qg-squares-theta}
\end{equation}
The same identity holds coefficientwise in $\IntegerNumbers[[q]]$.
\end{proposition}

\begin{proof}
In the $d$-fold Cauchy product, choosing $q^{x_j^2}$ from the $j$th factor
produces $q^{x_1^2+\cdots+x_d^2}$, and every ordered representation occurs
exactly once.  Analytically, absolute convergence permits the Cauchy product.
Formally, a fixed coefficient receives only finitely many contributions because
$x_j^2\le n$ whenever the tuple contributes to $q^n$.
\end{proof}

\begin{warning}[Generating functions are not coefficient evaluations]
The proposition identifies $r_d(n)$ as a theta-power coefficient, but by
itself gives no divisor-sum formula for that coefficient.  In particular, the
classical evaluations of $r_6(n)$ and $r_8(n)$ require additional input:
one proves modularity of the relevant theta powers, determines a basis (or a
dimension bound) for the appropriate modular-form space, decomposes into
Eisenstein and possible cusp components, and only then compares Fourier
coefficients.  Such evaluations must not be inferred from the Cauchy-product
argument alone.  The separate two- and four-square proofs below supply their
own additional arithmetic and analytic input.
\end{warning}

Define the real Dirichlet character modulo $4$ by
\[
 \chi_4(d)=
 \begin{cases}
  0,&2\mid d,\\
  1,&d\equiv1\pmod4,\\
 -1,&d\equiv3\pmod4.
 \end{cases}
\]

\begin{lemma}[Gaussian unique factorization]
\label{qg:lem-gaussian-ufd}
The ring $\IntegerNumbers[\ImaginaryUnit]$ is Euclidean for
$N(x+\ImaginaryUnit y)=x^2+y^2$, hence is a unique-factorization domain.
For every rational prime $p$:
\begin{enumerate}[label=\textup{(\roman*)}]
\item $2=-\ImaginaryUnit(1+\ImaginaryUnit)^2$;
\item if $p\equiv3\pmod4$, then $p$ is irreducible in
      $\IntegerNumbers[\ImaginaryUnit]$;
\item if $p\equiv1\pmod4$, then $p=\varpi\bar\varpi$ for a Gaussian prime
      $\varpi$, unique up to a unit and conjugation.  The alternate letter
      $\varpi$ prevents collision with the circular constant $\pi$ used in
      the theta and Fourier formulas above.
\end{enumerate}
\end{lemma}

\begin{proof}
For nonzero $\beta$, round the real and imaginary parts of $\alpha/\beta$
to obtain $\gamma\in\IntegerNumbers[\ImaginaryUnit]$.  Then
\[
 \AbsoluteValue{\alpha/\beta-\gamma}^2\le\frac12,
\]
and therefore
\[
 N(\alpha-\gamma\beta)
 =N(\beta)\AbsoluteValue{\alpha/\beta-\gamma}^2
 \le\frac12N(\beta)<N(\beta).
\]
Thus Euclidean division holds.  Descent in
the norm gives factorization into irreducibles, and B\'ezout's identity gives
Euclid's lemma, hence uniqueness up to order and the units
$\{\pm1,\pm\ImaginaryUnit\}$.

The factorization of $2$ is immediate.  If $p\equiv3\pmod4$ factored into
nonunits, the norms of the factors would force $p=x^2+y^2$, impossible
modulo $4$.  Thus $p$ is irreducible.

Now let $p\equiv1\pmod4$ and put
\[
 t=\left(\frac{p-1}{2}\right)!.
\]
Pairing $j$ with $p-j$ in Wilson's theorem gives
\[
 -1\equiv(p-1)!
 \equiv(-1)^{(p-1)/2}
 \left(\left(\frac{p-1}{2}\right)!\right)^2
 \equiv t^2\pmod p,
\]
because $(p-1)/2$ is even.  Hence
$p\mid(t+\ImaginaryUnit)(t-\ImaginaryUnit)$ in the Gaussian integers.
If $p$ were irreducible, it would be prime and would divide one factor, which
is impossible because its imaginary coordinate is $\pm1$.  Therefore
$p=\varpi\rho$ with nonunits.  Norms force
$N(\varpi)=N(\rho)=p$.  An element of prime norm is irreducible, so $\varpi$ and
$\rho$ are Gaussian primes.  Moreover,
\[
 \varpi\bar\varpi=N(\varpi)=p=\varpi\rho,
\]
and cancellation gives $\rho=\bar\varpi$.  If
$p=\sigma\bar\sigma$ is another such factorization, uniqueness of
factorization applied to $\varpi\bar\varpi=\sigma\bar\sigma$ shows that $\sigma$
is associated either to $\varpi$ or to $\bar\varpi$.  This proves uniqueness up to
a unit and conjugation.
\end{proof}

\begin{theorem}[Jacobi's two-square theorem]
\label{qg:thm-two-square}
For every $n\ge1$,
\begin{equation}
 r_2(n)=4\sum_{d\mid n}\chi_4(d).
\label{eq:qg-two-square}
\end{equation}
Equivalently, if
\[
 n=2^a\prod_{j=1}^{u}p_j^{\alpha_j}
       \prod_{\ell=1}^{v}r_\ell^{\beta_\ell},
 \qquad
 p_j\equiv1\pmod4,\quad r_\ell\equiv3\pmod4,
\]
then $r_2(n)=0$ if some $\beta_\ell$ is odd, and otherwise
$r_2(n)=4\prod_{j=1}^{u}(\alpha_j+1)$.
\end{theorem}

\begin{proof}
A representation $n=x^2+y^2$ is a Gaussian integer of norm $n$.  By
\Cref{qg:lem-gaussian-ufd}, the prime above $2$ has forced exponent up to a
unit.  A prime $r_\ell\equiv3\pmod4$ is self-conjugate and contributes an
even exponent to a norm, so $\beta_\ell$ must be even and then its exponent
in $x+\ImaginaryUnit y$ is forced.  For
$p_j=\varpi_j\bar\varpi_j$, the exponent of $\varpi_j$ may be any
$e_j\in\{0,\ldots,\alpha_j\}$ and the exponent of $\bar\varpi_j$ is then
$\alpha_j-e_j$.  Finally there are four units.  Unique factorization shows
that these choices are exhaustive and nonduplicative, proving the product
formula.

The divisor sum is multiplicative.  Its local factor is $1$ at $2$,
$\alpha_j+1$ at $p_j\equiv1\pmod4$, and
$1-1+\cdots+(-1)^{\beta_\ell}$ at $r_\ell\equiv3\pmod4$.  These are exactly
the factors in the product formula, proving \eqref{eq:qg-two-square}.
\end{proof}

\begin{corollary}[Two-square Lambert series]
\label{qg:cor-two-square-lambert}
For $\AbsoluteValue q<1$,
\begin{equation}
 \left(\sum_{m\in\IntegerNumbers}q^{m^2}\right)^2
 =1+4\sum_{d=1}^{\infty}\chi_4(d)\frac{q^d}{1-q^d}
 =1+4\sum_{j=0}^{\infty}(-1)^j
       \frac{q^{2j+1}}{1-q^{2j+1}}.
\label{eq:qg-two-square-lambert}
\end{equation}
\end{corollary}

\begin{proof}
By \Cref{qg:prop-squares-theta,qg:thm-two-square}, the coefficient of $q^n$
in the first Lambert series is $4\sum_{d\mid n}\chi_4(d)=r_2(n)$.
Since $\chi_4$ vanishes on even integers and equals $(-1)^j$ at $2j+1$,
the second form follows.
\end{proof}

\paragraph{Lean realization.}
The complete Gaussian-integer proof above is retained as the human proof.  The
formal proof in \texttt{JacobiTwoSquareCount} takes a complementary analytic
route: over $\ComplexNumbers$ it specializes Ramanujan's ${}_1\psi_1$
summation at base $Q=q^2$ with $a=-1$, $b=-Q$, and $z=q$, proves that the
bilateral summand is even, expands the resulting positive-index terms as an
absolutely summable double series, and reduces the product side by elementary
$q$-Pochhammer product identities.  Uniqueness of coefficients of convergent
power series then gives the arithmetic count.  Concretely,
\texttt{sumSqRep\_two\_eq\_four\_mul\_twoSquareDivisorSum} proves
\eqref{eq:qg-two-square} for every $n:\NaturalNumbers$ with $n\ne0$, as an
identity in $\IntegerNumbers$, while
\texttt{sumSqRep\_two\_eq\_four\_mul\_prod} gives the displayed product under
the explicit hypothesis that every prime divisor congruent to $3$ modulo $4$
has even valuation.  Finally,
\texttt{theta\_sq\_eq\_chi4\_lambert} and
\texttt{theta\_sq\_eq\_odd\_lambert} give the two respective Lambert forms
unconditionally for $q$ in any complete normed field with $\lVert q\rVert<1$;
thus the complex coefficient argument feeds the already general
hypothesis-parameterized analytic kernels rather than restricting the final
corollary to $\ComplexNumbers$.

\begin{theorem}[Jacobi's four-square theorem]
\label{qg:thm-four-square}
For every $n\ge1$,
\begin{equation}
 r_4(n)=8\sum_{\substack{d\mid n\\4\nmid d}}d.
\label{eq:qg-four-square}
\end{equation}
\end{theorem}

\begin{proof}
Fix $x\in\ComplexNumbers$ with $0<\AbsoluteValue x<1$.  Apply the
minus-sign form of \Cref{thm:jtp} with base $x$ and argument $a^2x$, and
multiply by $a$.  This gives
\begin{equation}
 (a-a^{-1})\prod_{\nu\ge1}
 (1-a^2x^\nu)(1-a^{-2}x^\nu)(1-x^\nu)
 =\sum_{m\in\IntegerNumbers}(-1)^ma^{2m+1}x^{m(m+1)/2}.
\label{eq:qg-four-square-start}
\end{equation}
For this fixed $x$, both sides and their $a$-derivatives converge normally on
compact annuli in the $a$-plane.  As functions of $x$, all theta series,
products, and logarithmic-derivative series used below converge locally
uniformly for $\AbsoluteValue x<1$, so the subsequent differentiations,
rearrangements, and parity dissections are legitimate.
Differentiate at $a=1$ and divide by $2$:
\begin{equation}
 (x;x)_\infty^3
 =\frac12\sum_{m\in\IntegerNumbers}(-1)^m(2m+1)x^{m(m+1)/2}.
\label{eq:qg-jacobi-cube}
\end{equation}
After squaring \eqref{eq:qg-jacobi-cube}, the summand indexed by $(m,n)$ is
\[
 \frac14(-1)^{m+n}(2m+1)(2n+1)
 x^{(m^2+n^2+m+n)/2}.
\]
If $m+n$ is even, put
\[
 m=r+s,\qquad n=r-s.
\]
Then
\[
 (2m+1)(2n+1)=(2r+1)^2-(2s)^2,
 \qquad
 \frac{m^2+n^2+m+n}{2}=r^2+s^2+r.
\]
If $m+n$ is odd, put
\[
 m=r+s,\qquad n=s-r-1.
\]
Now $(-1)^{m+n}=-1$, while
\[
 (2m+1)(2n+1)=(2s)^2-(2r+1)^2.
\]
Thus the parity sign reverses this coefficient, and the exponent is again
$r^2+s^2+r$.  Both changes of variables are bijections from their respective
parity classes to $\IntegerNumbers^2$, so the two classes contribute equally.
Consequently,
\begin{equation}
 (x;x)_\infty^6
 =\frac12\sum_{r,s\in\IntegerNumbers}
 \bigl((2r+1)^2-(2s)^2\bigr)x^{r^2+s^2+r}.
\label{eq:qg-four-square-coalesce}
\end{equation}

Let $D=x\,d/dx$ and set
\[
 A(x)=\sum_{s\in\IntegerNumbers}x^{s^2},
 \qquad
 B(x)=\sum_{r\in\IntegerNumbers}x^{r^2+r}.
\]
Then \eqref{eq:qg-four-square-coalesce} is
\begin{equation}
 (x;x)_\infty^6
 =\frac12\bigl[A(x)(1+4D)B(x)-4B(x)DA(x)\bigr].
\label{eq:qg-four-square-D}
\end{equation}
The triple product gives
\begin{align}
 A(x)&=\prod_{\nu\ge1}(1+x^{2\nu-1})^2(1-x^{2\nu}),
 \label{eq:qg-A-product}\\
 B(x)&=2\prod_{\nu\ge1}(1+x^{2\nu})^2(1-x^{2\nu}).
 \label{eq:qg-B-product}
\end{align}
Using
\[
 D\log(1+x^j)=\frac{jx^j}{1+x^j},
 \qquad
 D\log(1-x^j)=-\frac{jx^j}{1-x^j},
\]
the two copies of
$D\log\prod_{\nu\ge1}(1-x^{2\nu})$ cancel, and
\begin{align*}
 D\log B-D\log A
 &=2\sum_{\nu\ge1}\left(
   \frac{2\nu x^{2\nu}}{1+x^{2\nu}}
   -\frac{(2\nu-1)x^{2\nu-1}}{1+x^{2\nu-1}}
   \right).
\end{align*}
Since $A(x)B(x)/2$ is the product on the first line below,
\eqref{eq:qg-four-square-D} becomes
\begin{align}
 (x;x)_\infty^6
 &=\prod_{\nu\ge1}
 (1+x^{2\nu-1})^2(1+x^{2\nu})^2(1-x^{2\nu})^2\notag\\
 &\quad\times\left[
 1-8\sum_{\nu\ge1}\left(
 \frac{(2\nu-1)x^{2\nu-1}}{1+x^{2\nu-1}}
 -\frac{2\nu x^{2\nu}}{1+x^{2\nu}}
 \right)\right].
\label{eq:qg-four-square-logdiff}
\end{align}
We first verify the product identity
\[
 \prod_{\nu\ge1}(1+x^\nu)^4(1-x^\nu)^2
 =\prod_{\nu\ge1}
 (1+x^{2\nu-1})^2(1+x^{2\nu})^2(1-x^{2\nu})^2.
\]
Indeed, the quotient of the left product by the right product is
\[
 \left[
 \prod_{\nu\ge1}(1-x^{4\nu-2})
 \prod_{\nu\ge1}(1+x^{2\nu})
 \right]^2=1,
\]
because
\[
 \prod_{\nu\ge1}(1+x^{2\nu})
 =\frac{(x^4;x^4)_\infty}{(x^2;x^2)_\infty}
 =\frac1{(x^2;x^4)_\infty}
 =\prod_{\nu\ge1}(1-x^{4\nu-2})^{-1}.
\]
This identity therefore turns \eqref{eq:qg-four-square-logdiff} into
\begin{equation}
 \prod_{\nu\ge1}\left(\frac{1-x^\nu}{1+x^\nu}\right)^4
 =1-8\sum_{\nu\ge1}\left(
 \frac{(2\nu-1)x^{2\nu-1}}{1+x^{2\nu-1}}
 -\frac{2\nu x^{2\nu}}{1+x^{2\nu}}
 \right).
\label{eq:qg-four-square-quotient}
\end{equation}
The minus-sign triple product with base $x^2$ and argument $x$ gives
\[
 \sum_{m\in\IntegerNumbers}(-1)^mx^{m^2}
 =(x,x,x^2;x^2)_\infty.
\]
Furthermore,
\[
 (x,x,x^2;x^2)_\infty
 =(x;x^2)_\infty^2(x^2;x^2)_\infty
 =\frac{(x;x)_\infty}{(-x;x)_\infty}
 =\prod_{\nu\ge1}\frac{1-x^\nu}{1+x^\nu},
\]
where the middle equality follows from
\[
 (x;x)_\infty=(x;x^2)_\infty(x^2;x^2)_\infty,
 \qquad
 (x;x)_\infty(-x;x)_\infty=(x^2;x^2)_\infty.
\]

Replace $x$ by $-x$ in \eqref{eq:qg-four-square-quotient}.  Since
$m^2\equiv m\pmod2$, the preceding triple-product specialization turns the
left side into the fourth power of
$\sum_{m\in\IntegerNumbers}x^{m^2}$.  The right side becomes
\[
 1+8\left[
 \sum_{\substack{j\ge1\\j\ \mathrm{odd}}}
  \frac{jx^j}{1-x^j}
 +\sum_{\nu\ge1}\frac{2\nu x^{2\nu}}{1+x^{2\nu}}
 \right].
\]
For the even part,
\[
 \frac{2\nu x^{2\nu}}{1+x^{2\nu}}
 =\frac{2\nu x^{2\nu}}{1-x^{2\nu}}
  -\frac{4\nu x^{4\nu}}{1-x^{4\nu}},
\]
and hence
\[
 \sum_{\nu\ge1}\frac{2\nu x^{2\nu}}{1+x^{2\nu}}
 =\sum_{\substack{j\ge1\\2\mid j}}\frac{jx^j}{1-x^j}
  -\sum_{\substack{j\ge1\\4\mid j}}\frac{jx^j}{1-x^j}.
\]
Adding the odd-index sum leaves exactly the indices not divisible by $4$.
Thus
\begin{equation}
 \left(\sum_{m\in\IntegerNumbers}x^{m^2}\right)^4
 =1+8\sum_{\substack{j\ge1\\4\nmid j}}
 \frac{jx^j}{1-x^j}.
\label{eq:qg-four-square-lambert}
\end{equation}
By \Cref{qg:prop-squares-theta}, the coefficient of $x^n$ on the left is
$r_4(n)$; on the right it is
$8\sum_{d\mid n,\,4\nmid d}d$.  This proves
\eqref{eq:qg-four-square}.
\end{proof}

\begin{remark}
The two-square proof is arithmetic and the four-square proof analytic, but
both terminate in Lambert series.  This is the common mechanism: a theta
product exposes geometric denominators, and coefficient extraction turns
them into divisor sums.
\end{remark}
